\documentclass[11pt,a4paper]{article}
\usepackage[T1]{fontenc}
\usepackage{lmodern}
\usepackage[margin=28mm]{geometry}
\usepackage{amsmath,amssymb,amsthm,mathtools}
\usepackage{microtype}
\usepackage{enumitem}
\usepackage{xcolor}
\usepackage[colorlinks=true,linkcolor=blue!40!black,citecolor=blue!40!black,urlcolor=blue!40!black]{hyperref}
\setlength{\parskip}{3pt}
\setlength{\parindent}{1.2em}
\setlist{itemsep=2pt,topsep=4pt}
\allowdisplaybreaks[2]
\emergencystretch=2em
\newtheorem{theorem}{Theorem}[section]
\newtheorem{lemma}[theorem]{Lemma}
\newtheorem{proposition}[theorem]{Proposition}
\newtheorem{corollary}[theorem]{Corollary}
\theoremstyle{definition}
\newtheorem{definition}[theorem]{Definition}
\newtheorem{example}[theorem]{Example}
\theoremstyle{remark}
\newtheorem{remark}[theorem]{Remark}
\newcommand{\Q}{\mathbb Q}
\newcommand{\Z}{\mathbb Z}
\newcommand{\F}{\mathbb F}
\newcommand{\Pone}{\mathbb P^1}
\newcommand{\OK}{\mathcal O_K}
\newcommand{\Gm}{\mathbb G_m}
\newcommand{\cP}{\mathcal P}
\newcommand{\cL}{\mathcal L}
\newcommand{\cR}{\mathcal R}
\newcommand{\cB}{\mathcal B}
\newcommand{\tors}{\mathrm{tors}}
\DeclareMathOperator{\lcm}{lcm}
\DeclareMathOperator{\rank}{rank}
\DeclareMathOperator{\Gal}{Gal}
\DeclareMathOperator{\Spec}{Spec}
\hypersetup{pdftitle={Congruence Compression and Simultaneous Hilbert Irreducibility},pdfauthor={KEIO},pdfsubject={Thin sets, finitely generated multiplicative groups, and uniform specialization}}
\title{\bfseries Congruence Compression and\\Simultaneous Hilbert Irreducibility\\[5pt]\large With an independent application to Baker's question}
\author{KEIO\thanks{Research manuscript prepared with AI assistance. The arguments are supplied for independent mathematical review; no proof-assistant verification or historical priority certification is asserted.}}
\date{28 October 2026}
\begin{document}
\maketitle
\begin{abstract}
Let $K$ be a number field and let $T\subset K$ be thin. We prove that
its union of affine transforms $a(z)T+h(z)$, with $z$ in a fixed finitely
generated subgroup of a torus, has lattice density zero among algebraic
integers when $a$ is a nonzero monomial and $h$ is a Laurent polynomial.
In fixed integral-basis boxes of radius $H$, the exceptional count is
$O(H^{[K:\Q]}\exp\{-c\log H/\log\log H\})$. A separate congruence argument
gives the corresponding one-dimensional bound on the rational-integer
slice. The proof combines a prime-compression consequence of Pollack's
argument with a local deficit for each individual curve cover; the cover
index is kept fixed before the Chinese remainder theorem is applied.
This yields irreducibility of finitely many prescribed polynomial
specializations simultaneously under all allowed affine transforms, for
a density-one set of integral parameters. An elementary finite-field
value-set lemma also treats parameter-dependent polynomials of every
degree at least two whose leading coefficient is a monomial. The
classical inputs, the relation to an earlier square-value result, and
the limits of the conclusions are stated explicitly. An independent appendix proves rational linear independence of Dirichlet $L$-values at $1$ for a specified infinite family of moduli by a separate cyclotomic argument.
\end{abstract}

\noindent\textbf{Keywords.} Thin sets; Hilbert irreducibility; $S$-units;
finite-field value sets; congruence sieve; algebraic integers.

\section{Statements and conventions}

The purpose of this paper is to separate two sources of complexity in a
uniform specialization problem. A thin set has a local deficit at suitably
chosen primes, but an unbounded collection of multiplicative parameters
seems to introduce too many translates of that set. At a carefully chosen
product of primes, all parameters reduce through one small finite quotient.
The local deficits then dominate the number of parameter classes.

Fix a number field $K$ of degree $D=[K:\Q]$ and an integral basis
$\boldsymbol\omega=(\omega_1,\ldots,\omega_D)$ of $\OK$. For $H\ge1$, put
\[
 \cB_K(H)=\left\{\sum_{j=1}^D n_j\omega_j:
          n_j\in\Z,\ |n_j|\le H\right\}.
\]
Density below means density in these boxes. The conclusions hold for every
fixed choice of integral basis. All constants may depend on the fixed
arithmetic and geometric data. No uniformity in the degrees, field,
generators, or coefficients is implicit.

A Laurent polynomial in $z=(z_1,\ldots,z_k)$ belongs to
$K[z_1^{\pm1},\ldots,z_k^{\pm1}]$. A nonzero monomial means
\[
 a(z)=c z_1^{\alpha_1}\cdots z_k^{\alpha_k},
 \qquad c\in K^\times,\quad \alpha_j\in\Z.
\]
Let $\Gamma\le(K^\times)^k$ be finitely generated, with rank $r$ and finite
torsion subgroup of order $w$. For $N\ge1$, its subgroup of $N$th powers
is denoted by $\Gamma^N$; in particular,
\begin{equation}\label{eq:groupquotient}
 |\Gamma/\Gamma^N|\le wN^r.
\end{equation}

\begin{definition}\label{def:thin}
A subset $T\subset K$ is \emph{thin} if it is contained in a finite union
of a finite set and sets
\[
 f(C(K))\cap K,
\]
where $C$ is a smooth projective geometrically integral curve over $K$ and
$f:C\to\Pone_K$ is a morphism of degree at least two.
\end{definition}
This is the usual arithmetic notion of thinness on the affine line.
Normalization and projective completion give the displayed form from the
standard definition by generically finite covers, with finitely many
exceptional images absorbed into the finite set; see \cite[Chapter~3]{Serre}.
It is enough that $T$ be a subset of such a union.

\begin{theorem}[Affine transforms of a thin set]\label{thm:thin}
Let $T\subset K$ be thin, let $a$ be a nonzero monomial, and let $h$ be a
Laurent polynomial. Define
\begin{equation}\label{eq:Rthin}
 \cR=\OK\cap\{a(z)y+h(z):y\in T,\ z\in\Gamma\}.
\end{equation}
Then the following assertions hold.
\begin{enumerate}[label=\textup{(\roman*)}]
\item For every $\varepsilon>0$ there is a squarefree positive rational
integer $m$ such that
\[
 |\cR\bmod m\OK|\le\varepsilon m^D.
\]
\item There are $c>0$ and $H_0$ such that, for $H\ge H_0$,
\[
 |\cR\cap\cB_K(H)|\ll H^D
       \exp\left\{-c\frac{\log H}{\log\log H}\right\}.
\]
In particular, $\cR$ has density zero.
\item There is also a constant $c'>0$ such that
\[
 |\cR\cap\Z\cap[-H,H]|\ll H
       \exp\left\{-c'\frac{\log H}{\log\log H}\right\}.
\]
The last assertion concerns the rational-integer slice and is not inferred
from the ambient density statement.
\end{enumerate}
\end{theorem}

The theorem asserts density zero, not that the union in
\eqref{eq:Rthin} is itself thin. It allows arbitrarily large positive or
negative exponents in $\Gamma$ and arbitrary cancellation in $h(z)$.
The number of multiplicative coordinates and the Laurent polynomial are
fixed throughout.

\begin{theorem}[Simultaneous Hilbert irreducibility]\label{thm:HIT}
Let $P_1(T,Y),\ldots,P_s(T,Y)\in K(T)[Y]$ be irreducible, of positive
degree in $Y$, and let $G_i$ be the generic splitting-field Galois group
of $P_i$ over $K(T)$, with its action on the roots. Fix $a,h,\Gamma$ as
above. Outside a set of $\OK$ satisfying the bounds in
Theorem~\ref{thm:thin}(ii)--(iii), every $t\in\OK$ has the following
property: for every $1\le i\le s$ and every $z\in\Gamma$,
\[
 P_i\bigl(a(z)t+h(z),Y\bigr)
\]
is defined, has its generic degree, is separable and irreducible over
$K$, and has Galois group $G_i$ with the generic permutation action,
up to conjugacy.
\end{theorem}

The exceptional set controls the existential quantifier ``there is a bad
$z$''. Its complement therefore controls \emph{all} $z$ simultaneously.
This is different from a separate density statement for each fixed
translate. A finite collection of distinct pairs $(a,h)$ can also be
allowed, by taking a finite union of exceptional sets.

We will also prove a result that is not merely an application of a fixed
thin set, because its lower coefficients may vary with $z$.

\begin{theorem}[Polynomial families of arbitrary degree]\label{thm:polynomial}
Suppose $d\ge2$ and
\[
 F(X,z)=c z^{\boldsymbol\alpha}X^d+
                  \sum_{i=0}^{d-1} A_i(z)X^i,
 \qquad A_i\in K[z_1^{\pm1},\ldots,z_k^{\pm1}],\quad c\ne0.
\]
The set
\[
 \cR_F=\{t\in\OK:t=F(x,z)\text{ for some }x\in K,\ z\in\Gamma\}
\]
satisfies all three conclusions of Theorem~\ref{thm:thin}. There is no
parity restriction on $d$.
\end{theorem}

\subsection*{Relation to earlier work}
The use of moduli with small multiplicative exponent is classical.
Pollack \cite{Pollack} applies such moduli to sums from finitely generated
multiplicative groups over number fields. We use a stronger quantitative
form visible in his proof, not a new theorem about primes. The other
input is the standard local deficit for an individual thin cover; a
reference is Browning--Loughran \cite[Lemma~3.8(2)]{BL}. We give a direct
curve proof suited to the present application.

The transfer established here retains a free point on a curve while
allowing all multiplicative parameter values. An earlier corpus manuscript
\cite{Corpus60} treats an even-degree polynomial family over $\Q$ and
answers the square-value question raised after Example~1.3 of
Corvaja--Zannier \cite{CZ}. The present paper does not count that original
question as newly solved again. Its exact scope and the absence of a
historical priority claim are discussed in Section~\ref{sec:scope}.

\section{Primes with a compressed common exponent}

We isolate the prime-number input. Relative density refers to density
among rational primes.

\begin{proposition}[External sieve input]\label{prop:Pollack}
For every $\delta>0$ there is $\eta\in(0,1)$ such that, for all sufficiently
large $X$, the number of primes $p\le X$ for which $p-1$ has a prime
factor greater than $X^{1-\eta}$ is less than $\delta X/\log X$.
\end{proposition}
This is Pollack's Lemma~2.1 \cite{Pollack}, a consequence of the Brun-sieve
argument of Erd\H{o}s. It is the only sieve estimate assumed here. We also
use the prime number theorem and the Chebotarev density theorem.

\begin{lemma}[Prime compression]\label{lem:compression}
If $\cP$ has positive lower relative density, there are constants
$\delta,\eta>0$ such that, for every sufficiently large $X$, a finite
$\cL\subset\cP\cap[\delta X/8,X]$ satisfies
\begin{align}
 L:=|\cL|&\ge\frac{\delta X}{4\log X},\label{eq:Lbound}\\
 \log N&\le X^{1-\eta}\log X,
 &N&:=\lcm_{p\in\cL}(p-1),\label{eq:Nbound}\\
 \log m&\le2X,
 &m&:=\prod_{p\in\cL}p.\label{eq:mbound}
\end{align}
In particular, $\log N=o(L)$. If only positive upper relative density is
assumed, the same conclusions hold along an unbounded sequence of $X$.
\end{lemma}
\begin{proof}
Choose $\delta>0$ so that $|\cP\cap[1,X]|\ge\delta X/\log X$ for large
$X$. Proposition~\ref{prop:Pollack}, with parameter $\delta/2$, removes
at most $\delta X/(2\log X)$ primes. The prime number theorem removes at
most $\delta X/(4\log X)$ further primes below $\delta X/8$.
Take all remaining primes for $\cL$.
Every prime divisor of $N$ is at most $X^{1-\eta}$, and its largest
prime-power contribution is at most $X$. Thus
$N\le X^{X^{1-\eta}}$. Also
$\log m\le\sum_{p\le X}\log p\le2X$ for large $X$.
Finally,
\[
 \frac{\log N}{L}\le\frac4\delta X^{-\eta}(\log X)^2\longrightarrow0.
\]
For positive upper density, choose the initial sequence on which the
first counting inequality holds and repeat the argument.
\end{proof}

\begin{remark}
The separate estimates for $L$ and $\log N$ are essential. A bound of the
form $N<m^\kappa$ alone does not supply the comparison needed below.
All prime sets used in the applications are fixed completely split
Chebotarev sets with finitely many primes removed, so the lower-density
version applies for every sufficiently large $X$.
\end{remark}

\section{A local deficit for a curve cover}

We record exactly which standard geometric facts are used. A smooth
projective curve over a number field has a smooth proper model after
inverting finitely many primes. A fixed finite morphism and finitely
many auxiliary curves and maps extend after further exclusions;
geometric integrality and degree can be retained. A smooth projective
geometrically integral curve of genus $g$ over $\F_p$ satisfies
\begin{equation}\label{eq:Weil}
 \bigl|\#C(\F_p)-(p+1)\bigr|\le2g\sqrt p.
\end{equation}
We use these reduction facts and the Hasse--Weil bound as established
inputs, not as new results. The following argument is also consistent
with the general thin-set estimate in \cite[Lemma~3.8(2)]{BL}.

\begin{lemma}[Single-cover deficit]\label{lem:localgap}
Let $C/K$ be a smooth projective geometrically integral curve and let
$f:C\to\Pone_K$ have degree $d\ge2$. There are a finite Galois extension
$E/\Q$ containing $K$ and a finite set of rational primes $S$ such that,
whenever $p\notin S$ splits completely in $E$, for every prime ideal
$\mathfrak q$ of $K$ above $p$,
\begin{equation}\label{eq:localgap}
 \bigl|f_{\mathfrak q}(C_{\mathfrak q}(\F_p))\cap\F_p\bigr|
 \le\left(1-\frac1{2d}\right)p.
\end{equation}
The maps in this formula have good reduction and degree $d$.
\end{lemma}
\begin{proof}
Over an algebraic closure of $K$, the fiber product
$C\times_{\Pone}C$ contains the diagonal and an irreducible component
$D_0$ not equal to the diagonal. Indeed, in characteristic zero the
map is separable, and its generic fiber has $d$ distinct points; its
self-product contains ordered pairs with unequal entries. The
component $D_0$ and its normalization $\widetilde D_0$ are defined over
a finite extension of $K$. The normalization is a smooth projective
geometrically integral curve.

Take a finite Galois extension $E/\Q$ containing this field and all
conjugate data. Exclude finitely many rational primes so that the curves,
morphisms, degrees, and geometric integrality have good reduction.
The normalization map to $D_0\subset C\times C$ is an isomorphism away
from a fixed finite subset. The intersection of $D_0$ with the diagonal
is finite, because they are distinct components. After further prime
exclusions, the number of points removed for these reasons is uniformly
bounded in each residue field.

If $p$ splits completely in $E$, all these data reduce over $\F_p$ at
every prime of $E$ above $p$. The Hasse--Weil bound on
$\widetilde D_0$ therefore gives
\[
 p+O_f(\sqrt p)
\]
distinct ordered pairs $(P,Q)\in C_{\mathfrak q}(\F_p)^2$ with
$P\ne Q$ and $f(P)=f(Q)$. The error constant can be chosen uniformly
among the finitely many conjugate reductions. The same bound gives
$\#C_{\mathfrak q}(\F_p)=p+1+O_f(\sqrt p)$.

For $u\in\Pone(\F_p)$ let $r_u$ count the rational points of the fiber.
Since the reduced map has degree $d$, $0\le r_u\le d$. Consequently
\begin{align*}
 \#C_{\mathfrak q}(\F_p)-|f_{\mathfrak q}(C_{\mathfrak q}(\F_p))|
 &=\sum_{r_u>0}(r_u-1)\\
 &\ge\frac1d\sum_u r_u(r_u-1)\\
 &\ge\frac pd-O_f(\sqrt p).
\end{align*}
It follows that the projective image, and hence its finite part, has
cardinality at most $(1-1/d)p+O_f(\sqrt p)$. For all sufficiently large
$p$ this is at most $(1-1/(2d))p$. Absorb the remaining primes into $S$.
\end{proof}

\begin{lemma}[Reduction of a represented integral value]\label{lem:proper}
At a good prime ideal $\mathfrak q$, suppose $a(z)$ is a unit,
$h(z)$ is integral, and
\[
 t=a(z)f(P)+h(z)\in\OK,\qquad P\in C(K),\quad f(P)\in K.
\]
Then the reduction of $t$ lies in
\[
 \overline{a(z)}\,
   \bigl(f_{\mathfrak q}(C_{\mathfrak q}(\F_p))\cap\F_p\bigr)
       +\overline{h(z)}.
\]
\end{lemma}
\begin{proof}
The identity $f(P)=(t-h(z))/a(z)$ makes $f(P)$ integral at
$\mathfrak q$. Properness extends $P$ to a point of the proper model
over the valuation ring $\mathcal O_{K,\mathfrak q}$. The extended morphism
commutes with specialization. Since $f(P)$ is integral, its projective
specialization is finite and is exactly its residue in $\F_p$.
Reducing the displayed affine identity proves the assertion.
\end{proof}

No affine coordinate of $P$ has been assumed integral. This matters even
when $C=\Pone$: a rational point can have a denominator, and the map can
have poles. Projective good reduction and properness handle both issues.

\section{Compression and the Chinese remainder theorem}

Choose a finite presentation of $T$ as in Definition~\ref{def:thin}.
Treat each point of its finite part as a separate piece. Enlarge one
finite Galois extension $E/\Q$ to contain $K$ and the extensions needed
in Lemma~\ref{lem:localgap} for every cover. There are finitely many
pieces, say $J$; the empty case is immediate and may be discarded.

Choose a finite set $S$ of rational primes with the following properties.
Outside primes above $S$, every coordinate of every generator of $\Gamma$
is a unit; $c$ in $a(z)=cz^{\boldsymbol\alpha}$ is a unit; all coefficients
of $h$ and all finite exceptional values are integral; and every curve
and map used above has good reduction. Since only finitely many elements
of $K$ are involved, such $S$ exists. It follows that these unit and
integrality assertions hold for \emph{every} $z\in\Gamma$, including
negative powers of generators.

Let $\cP$ consist of rational primes outside $S$ that split completely in
$E$. Chebotarev gives positive density, and
\begin{equation}\label{eq:split}
 \OK/p\OK\simeq\prod_{\mathfrak q\mid p}\F_p\simeq\F_p^D
 \qquad(p\in\cP).
\end{equation}
After removing finitely many additional primes, there is one
$\theta\in(0,1)$ such that each piece has at most $\theta p$ possible
finite residues at each $\mathfrak q\mid p$. A singleton has at most one
residue, so it satisfies the same assertion for large $p$.

\begin{proposition}[Finite-modulus bound]\label{prop:CRT}
For any finite $\cL\subset\cP$, with
$m=\prod_{p\in\cL}p$, $N=\lcm_{p\in\cL}(p-1)$, and $L=|\cL|$,
\begin{align}
 \frac{|\cR\bmod m\OK|}{m^D}
 &\le JwN^r\theta^{DL},\label{eq:CRT}\\
 \frac{|(\cR\cap\Z)\bmod m\Z|}{m}
 &\le JwN^r\theta^L.\label{eq:CRTslice}
\end{align}
\end{proposition}
\begin{proof}
For each $p\in\cL$ and $\mathfrak q\mid p$, the multiplicative group
$\F_p^\times$ has exponent $p-1$, which divides $N$. Therefore the
simultaneous reduction of $\Gamma$ at all these primes factors through
$\Gamma/\Gamma^N$. Its size is at most $wN^r$ by
\eqref{eq:groupquotient}.

Fix \emph{both} a piece $i$ of the presentation of $T$ and a class
$\xi\in\Gamma/\Gamma^N$. At every $\mathfrak q\mid p\in\cL$, the
reductions of $a(z)$ and $h(z)$ are fixed for $z\in\xi$, and the former
is nonzero. The affine map $y\mapsto\overline{a(z)}y+\overline{h(z)}$
is a bijection of $\F_p$. Lemma~\ref{lem:proper}, or direct reduction for
a singleton piece, therefore places the corresponding values of $t$ in
at most $\theta p$ residue classes at that prime ideal.

The Chinese remainder theorem and \eqref{eq:split} show that, for this
fixed pair $(i,\xi)$, there are at most
\[
 \prod_{p\in\cL}(\theta p)^D=\theta^{DL}m^D
\]
possible residue classes modulo $m\OK$. Only now take the union over
$i$ and $\xi$. This proves \eqref{eq:CRT}; no independence of the
parameter residues at distinct primes is needed.

For rational integers, choose one prime ideal $\mathfrak q_p$ over
each $p\in\cL$. The residue of an integer in $\OK/\mathfrak q_p$
is its usual residue in $\F_p$. Repeating the argument with just these
$L$ coordinates gives at most $\theta^Lm$ residues for a fixed
$(i,\xi)$ and proves \eqref{eq:CRTslice} after taking the same union.
\end{proof}

\begin{remark}[Why the finite union is taken last]\label{rem:union}
One cannot assume a local deficit for the union of all thin pieces.
For example, the three thin sets defined by the existence of a rational
solution to
\[
 y^2=t,\qquad y^2=t-1,\qquad y^2=t(t-1)
\]
cover all finite-field residues at every odd prime: if neither $t$ nor
$t-1$ is a square, their product is a square. Their union is not all of
$\Q$, since $t=3$ gives the three nonsquares $3,2,6$. In
Proposition~\ref{prop:CRT}, membership in one fixed global piece is
imposed across every prime before the finite union is taken.
\end{remark}

\begin{proof}[Proof of Theorem~\ref{thm:thin}]
Apply Lemma~\ref{lem:compression} to $\cP$. Since
$\log N=o(L)$, the logarithms of the right sides of
\eqref{eq:CRT} and \eqref{eq:CRTslice} tend to $-\infty$. The first
residue assertion follows. Each residue class modulo $m\OK$ has density
$m^{-D}$ in the chosen boxes, so the same assertion already proves
qualitative density zero.

For the quantitative conclusion, take
\[
 X=\tfrac14\log H
\]
for sufficiently large $H$. The compression construction gives
$m\le e^{2X}\le H^{1/2}$. Write $\beta=-\log\theta>0$.
By \eqref{eq:Lbound}--\eqref{eq:Nbound},
\[
 \log(Jw)+r\log N=o(X/\log X),
 \qquad L\ge\frac{\delta X}{4\log X}.
\]
Thus, for all large $H$, the residue proportions in
\eqref{eq:CRT} and \eqref{eq:CRTslice} are respectively bounded by
\[
 \exp\{-D\beta\delta X/(8\log X)\},
 \qquad \exp\{-\beta\delta X/(8\log X)\}.
\]
Each residue class modulo $m\OK$ contains at most
$(2H/m+2)^D$ elements of $\cB_K(H)$. Multiplying this by the number of
admissible classes bounds the exceptional count by
\[
 (2H+2m)^D\exp\{-D\beta\delta X/(8\log X)\}.
\]
Since $m\le H^{1/2}$ and $X=\frac14\log H$, this has the required form.
For the rational-integer slice, each class modulo $m$ contains at most
$2H/m+2$ integers in $[-H,H]$. Use \eqref{eq:CRTslice} in exactly this
one-dimensional count. This proves (iii) separately.
\end{proof}

\section{Simultaneous irreducibility and examples}

We spell out the standard thin-set specialization argument, including
full Galois groups. The specialization of a finite Galois extension at
an unramified parameter has group a subgroup of the generic group;
quotient covers detect its fixed cosets. These are classical facts of
specialization theory; see \cite[Proposition~3.3.1]{Serre}.

\begin{lemma}[Thin bad-specialization set]\label{lem:hilbertbad}
For the finite list $P_i$ in Theorem~\ref{thm:HIT}, the parameters at
which any of the stated degree, separability, irreducibility, or Galois
group properties fails are contained in a thin subset of $K$.
\end{lemma}
\begin{proof}
Fix $i$ and let $L_i/K(T)$ be its splitting field. Remove the finite set
of poles, leading-coefficient zeros, discriminant zeros, and branch
values. At another parameter $y\in K$, its specialized Galois group
$G_{i,y}$ is a subgroup of $G_i$. If it is proper, it is contained in a
conjugate of a maximal proper subgroup $H<G_i$. The specialized action
on $G_i/H$ then fixes a coset. Equivalently, $y$ lifts to a $K$-point on
the smooth projective quotient curve with function field $L_i^H$.
Its map to $\Pone$ has degree $[G_i:H]\ge2$.

There are finitely many such subgroups. A smooth integral quotient curve
that is not geometrically integral has a constant field strictly larger
than $K$ and has no $K$-point; it contributes no values. The geometrically
integral quotient curves give pieces of Definition~\ref{def:thin}, and
the finitely many removed values give its finite part. Thus parameters
with a proper specialized group form a thin set. With the full generic
group, its action remains transitive because $P_i$ is generically
irreducible, so the specialized polynomial is irreducible as well.
Take the finite union over $i$.
\end{proof}

\begin{proof}[Proof of Theorem~\ref{thm:HIT}]
Let $T$ be the thin bad set from Lemma~\ref{lem:hilbertbad}. A failing
$t$ has, for some $z\in\Gamma$,
\[
 y=a(z)t+h(z)\in T,
 \qquad t=a(z)^{-1}y-a(z)^{-1}h(z).
\]
The reciprocal $a^{-1}$ is a nonzero monomial and $-a^{-1}h$ is a
Laurent polynomial. Theorem~\ref{thm:thin}, applied to these two
functions, bounds the set of all failing $t$ and proves both the
ambient and rational-integer conclusions.
\end{proof}

The inverse-affine formulation $P_i(a(z)^{-1}(t-h(z)),Y)$ is equivalent,
by the same substitution. One may also preserve the full generic Galois
group of the compositum of all splitting fields: apply the maximal
subgroup argument to that one extension. This retains generic joint
Galois relations, not merely the individual groups.

Universal Hilbert subsets have a substantial earlier literature.
D\`ebes--Zannier \cite{DZ} construct explicit universal sequences, and
Zannier's overview \cite[Theorem~H]{Zannier} describes further power-sum
results. Those statements fix a sequence and, for each fixed polynomial,
allow finitely many bad sequence members. Here the polynomials and
multiplicative affine family are fixed, and a density-one set of
translations is good for \emph{every} member of that family. This
quantifier distinction is useful, but is not a claim to the first
universal Hilbert theorem. In particular, no one parameter is asserted
to work for all polynomials chosen afterwards.

\begin{corollary}[Fixed numbers of multiplicative translates]\label{cor:sum}
Let $G\le K^\times$ be finitely generated, let $T\subset K$ be thin,
and let $b_1,\ldots,b_s\in K$ be fixed. Then
\[
 \OK\cap(T+b_1G+\cdots+b_sG)
\]
satisfies the three conclusions of Theorem~\ref{thm:thin}. The same is
true with $T$ replaced by $GT$.
\end{corollary}
\begin{proof}
Use $\Gamma=G^s$, $a=1$, and $h(z)=\sum b_jz_j$. For the final
assertion use $\Gamma=G^{s+1}$ and $a(z)=z_0$.
\end{proof}

\begin{example}[Odd powers and signed shifts]\label{ex:odd}
For every $d\ge2$, the image $\{x^d:x\in K\}$ is thin. For fixed
finitely many $b_j$ and $G$, almost every algebraic integer $t$ admits
no representation
\[
 t+b_1u_1+\cdots+b_su_s=x^d,
 \qquad x\in K,\quad u_j\in G.
\]
The same density-one statement holds among rational integers, even when
$K\ne\Q$. In particular, $d=3$ is included. Negative exponents in the
$u_j$ are unrestricted, so a positive-term counting argument would not
be an adequate substitute for the congruence proof.
\end{example}

\begin{example}[A nonpolynomial rational map]
The map $f(x)=x+x^{-1}$ has degree two as a map $\Pone\to\Pone$.
Its finite values form a thin set. Thus almost every $t$ avoids
\[
 t+b_1u_1+\cdots+b_su_s=x+x^{-1},
 \qquad x\in K^\times,\quad u_j\in G,
\]
with the quantitative exceptional bounds above. Points at a pole are
excluded by the requirement that the displayed value be finite.
\end{example}

\begin{example}[A cover from a nonrational curve]
Let $C/\Q$ be the smooth projective completion of
\[
 v^2=u^3-2.
\]
It is a genus-one curve, hence not a rational curve. The coordinate
$u$ defines a degree-two morphism $f:C\to\Pone$. Its finite rational
image is thin. Consequently, for fixed rational coefficients $b_j$ and
a fixed finitely generated $G\le\Q^\times$, the integers of the form
\[
 u+b_1u_1+\cdots+b_su_s,
 \qquad (u,v)\in C(\Q),\quad u_j\in G,
\]
have the asserted density-zero bound. The theorem applies whether this
particular curve has finitely or infinitely many rational points; no
rank computation is needed. This example illustrates that the proof is
not restricted to rational parametrizations.
\end{example}

\begin{example}[Uniform irreducibility]
The polynomial $Y^3-Y-T$ is irreducible in $\Q(T)[Y]$: the extension
$\Q(Y)/\Q(Y^3-Y)$ has degree three. Fix a finitely generated
$G\le\Q^\times$ and rational $b_1,\ldots,b_s$. Outside a quantitatively
sparse set of integers $t$, the polynomial
\[
 Y^3-Y-\bigl(u_0t+b_1u_1+\cdots+b_su_s\bigr)
\]
is irreducible over $\Q$ for every $(u_0,\ldots,u_s)\in G^{s+1}$.
The conclusion can be required simultaneously for any fixed finite
list of irreducible families.
\end{example}

\section{Polynomial families with varying coefficients}

Theorem~\ref{thm:polynomial} allows the lower coefficients to vary with
$z$, so the fixed-cover proof cannot simply be reused. A direct
finite-field argument gives the needed uniform gap.

\begin{lemma}[Elementary value-set bound]\label{lem:values}
Let $p$ be prime, let $d\ge2$ divide $p-1$, and let
$G\in\F_p[X]$ have degree exactly $d$. Then
\[
 |G(\F_p)|\le p-\frac{p-1}{d}.
\]
\end{lemma}
\begin{proof}
Put $k=(p-1)/d$ and let $a\ne0$ be the leading coefficient. The sums
of powers in $\F_p$ imply
\[
 S_j:=\sum_{x\in\F_p}G(x)^j=0\quad(1\le j<k),
 \qquad S_k=-a^k\ne0.
\]
Indeed, below degree $p-1$ every nonconstant monomial has zero sum,
as does the constant monomial in characteristic $p$; at degree $p-1$
only the leading term contributes.

Set $Q(T)=\prod_{x\in\F_p}(T-G(x))$. Since $k<p$, Newton's identities
with invertible indices give
\[
 Q(T)=T^p+O(T^{p-k}),
\]
where the notation means that all remaining terms have degree at most
$p-k$. Let $V=G(\F_p)$, let $m=p-|V|$, and let $r_v$ be the multiplicity
of $v$ in this list of values. Define monic degree-$m$ polynomials
\[
 A(T)=\prod_{v\in V}(T-v)^{r_v-1},
 \qquad B(T)=\prod_{b\in\F_p\setminus V}(T-b).
\]
They satisfy
\begin{equation}\label{eq:AB}
 Q(T)B(T)=(T^p-T)A(T).
\end{equation}
If $m<k$, then $p-k+m\le p-1$ and $m+1<p$. The terms of degrees
at least $p$ on the left of \eqref{eq:AB} are exactly $T^pB(T)$; on the
right they are exactly $T^pA(T)$. Thus $A=B$, and cancellation gives
$Q=T^p-T$. This would make $G$ a permutation of $\F_p$. But then
$S_k=\sum_{y\in\F_p}y^k=0$, because $1\le k<p-1$, a contradiction.
Therefore $m\ge k$.
\end{proof}

\begin{remark}
The lemma is a special finite-field value-set estimate, not a new
classification of exceptional polynomials. The displayed proof is
included so that neither a parity assumption nor such a classification
is needed in the application.
\end{remark}

\begin{proof}[Proof of Theorem~\ref{thm:polynomial}]
Choose a finite exceptional set $S$ so that, outside primes above $S$,
all generator coordinates are units, $c$ is a unit, and every
coefficient of every $A_i$ is integral. Take a finite Galois extension
$E/\Q$ containing $K$ and $\Q(\zeta_d)$. The rational primes outside
$S$ splitting completely in $E$ have positive density, split completely
in $K$, and satisfy $p\equiv1\pmod d$.

At a prime ideal $\mathfrak q$ over one of these primes, fix $z\in\Gamma$.
The specialized polynomial $F(X,z)$ has integral coefficients and unit
leading coefficient. If $v_{\mathfrak q}(x)<0$, its leading term has
valuation $d v_{\mathfrak q}(x)$, strictly smaller than the valuation of
every lower-degree term: for $i<d$,
\[
 v_{\mathfrak q}(A_i(z)x^i)\ge i v_{\mathfrak q}(x)
       >d v_{\mathfrak q}(x).
\]
A sum with a unique term of least valuation has that valuation. Therefore
$F(x,z)$ could not be integral. Every represented integral value thus
has $x$ integral at $\mathfrak q$.

The reduction of $F(X,z)$ has degree exactly $d$.
Lemma~\ref{lem:values} gives a relative value-set bound
\[
 1-\frac{p-1}{dp}\le\theta:=1-\frac1{2d}<1.
\]
For any compressed set of primes, fix a class
$\xi\in\Gamma/\Gamma^N$. Its reductions fix all the coefficients of
$F(X,z)$ simultaneously, so the bound holds for that one polynomial at
each prime ideal. The same CRT calculation as
Proposition~\ref{prop:CRT}, now with $J=1$, yields
\[
 \frac{|\cR_F\bmod m\OK|}{m^D}\le wN^r\theta^{DL}.
\]
Selecting one prime ideal above each rational prime similarly gives
$wN^r\theta^L$ for the rational-integer slice. Apply
Lemma~\ref{lem:compression} and the box count from the proof of
Theorem~\ref{thm:thin}. All three assertions follow.
\end{proof}

\begin{example}[A genuinely parameter-dependent odd-degree family]
Let $\Gamma\le(K^\times)^2$ be finitely generated. The integral values of
\[
 F(X,z_1,z_2)=z_1X^3+(z_2+z_2^{-1})X^2+(1+z_1^{-1})X+z_2
\]
with $X\in K$ and $(z_1,z_2)\in\Gamma$ form a set of density zero,
with both quantitative bounds. The quadratic and linear coefficients
vary, so this assertion is not obtained by identifying one fixed image
$f(C(K))$ and translating it.
\end{example}

\section{Fixed degrees and unrestricted proper powers}\label{sec:degrees}

The finite-union principle allows any fixed finite collection of
proper-power degrees. It does not allow all $d\ge2$ simultaneously.
This is not just a matter of losing a uniform constant in the proof.

\begin{proposition}[Squarefree-modulus obstruction]\label{prop:allpowers}
Let $U=\{x^d:x\in\Q,\ d\ge2\}$. For every squarefree positive integer
$M$, the integer proper powers, and hence $U\cap\Z$, map surjectively
onto $\Z/M\Z$.
\end{proposition}
\begin{proof}
For $M=1$ the assertion is immediate. Otherwise choose a prime $q$
not dividing $\lcm_{p\mid M}(p-1)$. The map $x\mapsto x^q$ is a
bijection of each $\F_p^\times$ and fixes zero, hence is a bijection
of $\F_p$. The Chinese remainder theorem makes it a bijection of
$\Z/M\Z$. For each residue, an integer representative of its preimage
has an integer $q$th power in that residue.
\end{proof}

Thus the arbitrarily small squarefree-residue conclusion of
Theorem~\ref{thm:thin} is false for $U$ itself. A countable union of
fixed-degree conclusions cannot prove it. Although integer proper
powers without shifts have density zero by elementary counting, the
present theorem does not prove the corresponding assertion after all
allowed $S$-unit shifts. Nor can the exponent be bounded merely from
the target size: the identity
\[
 -1+2^q+3^0=2^q
\]
gives a fixed target with representations of unbounded prime degree.
An all-degree tail estimate would require an additional argument that
handles such cancellation. None is claimed here.

\section{Scope and remaining limitations}\label{sec:scope}

The main conclusion is a transfer theorem: thin images and suitable
polynomial value sets remain sparse among integral targets after all
allowed multiplicative affine operations. It does not classify the
solutions for a fixed target, prove that the resulting set is thin,
or settle a general local--global principle.

The monomial assumption has a specific role. It makes $a(z)$, or the
leading coefficient of $F(X,z)$, a unit at every chosen prime for every
$z\in\Gamma$. A general Laurent polynomial can vanish modulo selected
primes even if it never vanishes on $\Gamma$ in characteristic zero.
The proof does not justify replacing the monomial by an arbitrary
nonvanishing function.

Likewise, a general parameter-dependent rational map $F(X,z)$ need not
have one fixed extension over which all its relevant off-diagonal
components are defined. The fixed Chebotarev set in the proof may then
be unavailable. No theorem for all such rational families is claimed.
The number of multiplicative summands is fixed, and no assertion is
made for an unbounded number of summands. The quantitative saving is
subpower, not a fixed power saving.

The square-value question stated by Corvaja and Zannier concerns whether
many integers $t$ avoid all squares $t+a2^m+b3^n$, with $m,n\in\Z$,
and analogous $S$-unit shifts \cite[after Example~1.3]{CZ}. It is a
legitimate earlier question, and the square specialization here answers
it. However, that conclusion already occurs in \cite{Corpus60}; it is
not a new resolution counted in this paper. The extensions to arbitrary
thin covers, simultaneous affine specialization, the separate
rational-integer slice, and all-degree number-field polynomial families
are the results proved here. Their historical priority has not been
established by an exhaustive literature search.

This manuscript does not establish a new complete solution of a previously
named open problem under all the stated novelty requirements. It does not resolve Baker's question for arbitrary moduli. Appendix~\ref{bkapp:section}
proves only its displayed family by an independent method. No assertion
is made on the rational-base Lambert conjecture at $3/2$ or on Mills-type
constants. The proofs are unconditional relative to the explicitly
identified classical theorems, but have not been certified in a proof
assistant or subjected to external peer review.

\clearpage
\appendix
\section{An independent application to Baker's question}
\label{bkapp:section}

Baker asked whether, for every integer $q\geq3$, the numbers
$L(1,\chi)$ as $\chi$ ranges over the nonprincipal Dirichlet characters
modulo $q$ are linearly independent over $\mathbb Q$
\cite[p.~48]{BK_Baker1975}; see \cite{BK_KR2022} for the history and
the precise parity reduction.  The coprimality case
$\gcd(q,\varphi(q))=1$ follows from Baker--Birch--Wirsing
\cite{BK_BBW1973}.
The argument below is independent of the thin-set construction in the main
text.  It proves an unconditional infinite family by a signed Galois
interchange, not by a general statement about arbitrary twisted cotangent
sums.  It does not resolve the question for all moduli or all prime powers,
and no claim of bibliographic priority is made.

Write $\lambda(n)$ for the exponent of
$(\mathbb Z/n\mathbb Z)^\times$, with $\lambda(1)=1$.

\begin{theorem}
\label{bkapp:main}
Let $m\geq1$ be odd and suppose that
\[
  \gcd(m,\lambda(m))=1,
  \qquad 8\nmid\lambda(m).
\]
Then the numbers $L(1,\chi)$, with $\chi$ ranging over all nonprincipal
Dirichlet characters modulo $32m$, are linearly independent over
$\mathbb Q$.  In particular, this holds for $q=32r$ whenever $r$ is an
odd prime satisfying $r\not\equiv1\pmod8$.
\end{theorem}

The final assertion supplies infinitely many moduli, since the primes
$r\equiv3\pmod8$ alone form an infinite set.  For these moduli
$\gcd(q,\lambda(q))=8$; thus they are not covered by the elementary
singleton-fibre criterion proved below.

\subsection{The external parity input}

We use two established facts for a fixed modulus $q>2$:
\begin{enumerate}
\item The values $L(1,\chi)$ at nonprincipal even characters are linearly
independent over $\overline{\mathbb Q}$, by Murty--Murty
\cite{BK_MM2011}; the statement is also recorded explicitly in
\cite[Theorem~3]{BK_KR2022}.
\item The $\mathbb Q$-span of these even values has zero intersection
with the $\mathbb Q$-span of the odd values
\cite[Theorem~5]{BK_KR2022}.
\end{enumerate}
The theorem numbering here refers to the cited arXiv version of
\cite{BK_KR2022}.
For the second fact, the underlying input is Baker's theorem applied to
logarithms of positive real algebraic numbers together with
$\log(-1)=\pi i$; see also the logarithm lemma of Murty--Saradha
\cite{BK_MS2010}, reproduced as Lemma~4 in \cite{BK_KR2022}.
An even value is an algebraic linear combination of such real logarithms,
whereas an odd value is an algebraic multiple of $\pi$.
Consequently it suffices to prove rational linear independence of the
odd values.  The restriction-fibre framework is developed in the earlier corpus
manuscript \cite{Corpus53} and is reproved below for completeness.
The additional step here is the signed interchange within the
two-element fibres and its application to the displayed modulus family.

\subsection{Algebraic representatives and restriction fibres}

For $q\geq3$, put
\[
 U_q=(\mathbb Z/q\mathbb Z)^\times,
 \qquad \lambda=\lambda(q),\qquad
 N=\operatorname{lcm}(q,\lambda),\qquad
 \zeta_n=e^{2\pi i/n}.
\]
For $a\in U_q$ and an odd character $\chi$ modulo $q$, define
\[
 y_a=\frac{\zeta_q^a+1}{\zeta_q^a-1}
     =-i\cot\frac{\pi a}{q},
 \qquad
 W_\chi=\sum_{a\in U_q}\chi(a)y_a
       \in\mathbb Q(\zeta_N).
\]
The partial-fraction expansion of the cotangent gives
\begin{equation}
\label{bkapp:L-W}
 L(1,\chi)=\frac{\pi}{2q}
   \sum_{a\in U_q}\chi(a)\cot\frac{\pi a}{q}
 =\frac{\pi i}{2q}W_\chi.
\end{equation}
Indeed, oddness makes the symmetric sum of $\chi(n)/n$ over nonzero
integers equal to twice the ordinary Dirichlet series; grouping the
symmetric sum by residue classes and using
$\sum_{k\in\mathbb Z}(a+kq)^{-1}=(\pi/q)\cot(\pi a/q)$ proves
\eqref{bkapp:L-W}.  All sums over $k$ here are symmetric limits.
Dirichlet nonvanishing implies $W_\chi\ne0$.

For $u\in U_N$, let $\sigma_u$ denote the automorphism of
$\mathbb Q(\zeta_N)$ sending $\zeta_N$ to $\zeta_N^u$.
Since character values have orders dividing $\lambda$, a change of
variable $b=ua$ gives
\begin{equation}
\label{bkapp:galois}
 \sigma_u(W_\chi)
 =\sum_{a\in U_q}\chi(a)^u y_{ua}
 =\overline{\chi^u}(u)W_{\chi^u}.
\end{equation}
The exponent $u$ on a character and its argument $u$ are interpreted
modulo the character order and modulo $q$, respectively.
Only the algebraic numbers $W_\chi$ are subjected to Galois
automorphisms; no action on the transcendental common factor in
\eqref{bkapp:L-W} is asserted.

\begin{lemma}[Restriction-fibre splitting]
\label{bkapp:fibres}
Put
\[
 \delta=\gcd(q,\lambda),\qquad
 K=\mathbb Q(\zeta_\lambda),\qquad
 H=\ker(U_q\longrightarrow U_\delta).
\]
Every rational linear relation among the odd numbers $W_\chi$ splits
into separate relations among characters having the same restriction
to $H$.  Each nonempty fibre is a coset of the group of even
characters modulo $\delta$, viewed as characters modulo $q$.
In particular, if $\delta=4$ all odd $W_\chi$ are linearly independent
over $\mathbb Q$.  If $\delta=8$, every odd fibre is a pair
$\{\chi,\chi\eta_8\}$, where
\[
 \eta_8(a)=(-1)^{(a^2-1)/8}\qquad(a\text{ odd})
\]
is the nonprincipal even character modulo $8$.
\end{lemma}

\begin{proof}
The reduction map $U_q\longrightarrow U_\delta$ is onto: a unit
modulo a divisor of $q$ lifts to a unit modulo $q$ by the Chinese
remainder theorem at the prime-power factors.
For each $s\in H$, the congruences
\[
 u_s\equiv s\pmod q,
 \qquad u_s\equiv1\pmod\lambda
\]
have a solution, because $s\equiv1\pmod\delta$.
The solution is a unit modulo $N$ and is unique modulo $N$.
These automorphisms identify $H$ with
$\operatorname{Gal}(\mathbb Q(\zeta_N)/K)$.
Equation~\eqref{bkapp:galois} then reads
\[
 \sigma_{u_s}(W_\chi)=\overline\chi(s)W_\chi.
\]
All characters of $H$ have values in $K$, since the exponent of $H$
divides $\lambda$.  For a character $\mu$ of $H$, the $K$-linear
operator
\[
 E_\mu=\frac1{|H|}\sum_{s\in H}\mu(s)\sigma_{u_s}
\]
sends $W_\chi$ to $W_\chi$ if $\chi|_H=\mu$, and to zero otherwise,
by character orthogonality.  Applying $E_\mu$ to a rational relation
proves the splitting assertion.

Two characters agree on $H$ exactly when their quotient factors
through $U_q/H=U_\delta$.  If both are odd, their quotient is even;
conversely multiplying by an even character modulo $\delta$
preserves both oddness and the restriction to $H$.
Modulo $4$ there is no nonprincipal even character, so every fibre
is a singleton and nonvanishing proves independence.  Modulo $8$
the even character group is $\{1,\eta_8\}$, which proves the last
assertion.
\end{proof}

\subsection{A signed Galois interchange}

\begin{lemma}
\label{bkapp:negative-swap}
Let $\eta$ be an even quadratic character modulo $q$, and let $\chi$
be odd.  Suppose that a unit $u\in U_N$ satisfies
\[
 \chi^u=\chi\eta,
 \qquad \eta(u)=-1.
\]
Then $W_\chi$ and $W_{\chi\eta}$ are linearly independent over
$\mathbb Q$.  In particular,
$L(1,\chi\eta)/L(1,\chi)$ is algebraic and irrational.
\end{lemma}

\begin{proof}
All these values are nonzero.  A rational dependence would therefore
give
\[
 c=\frac{W_{\chi\eta}}{W_\chi}\in\mathbb Q^\times.
\]
The integer $u$ is odd because $\lambda$ is even, and hence
$\eta^u=\eta$.  Thus $(\chi\eta)^u=\chi$.
By \eqref{bkapp:galois},
\[
 \sigma_u(c)
 =\frac{\overline\chi(u)W_\chi}
        {\overline{\chi\eta}(u)W_{\chi\eta}}
 =-\frac1c.
\]
Rationality would give $c=\sigma_u(c)$ and then $c^2=-1$, a
contradiction.  Formula~\eqref{bkapp:L-W} identifies the two
ratios and also proves that the ratio of the $L$-values is algebraic.
\end{proof}

\begin{proof}[Proof of Theorem~\ref{bkapp:main}]
Set $\ell=\lambda(m)$ and $q=32m$.
The hypotheses imply
\[
 \lambda(q)=\operatorname{lcm}(8,\ell),
 \qquad \gcd(q,\lambda(q))=8.
\]
Indeed the odd part of the gcd is trivial because
$\gcd(m,\ell)=1$, while the $2$-part of $\lambda(q)$ is exactly
$8$ because $8\nmid\ell$.
Lemma~\ref{bkapp:fibres} therefore reduces the odd-value assertion
to the independence of each pair $\{W_\chi,W_{\chi\eta_8}\}$.

By the Chinese remainder theorem, write
$\chi=\alpha\beta$ with $\alpha$ modulo $32$ and $\beta$ modulo
$m$.  The convention for $m=1$ is the trivial character.
Since
\[
 U_{32}=\langle-1\rangle\times\langle5\rangle,
 \qquad \operatorname{ord}_{32}(5)=8,
\]
we may write $\alpha(5)=\zeta_8^j$ for $0\leq j<8$.

First suppose that $j$ is odd.  Choose an integer $u$ satisfying
\begin{equation}
\label{bkapp:CRT}
 u\equiv5\pmod{32},
 \qquad u\equiv1\pmod{m\ell}.
\end{equation}
These congruences are compatible because
\[
 \gcd(32,m\ell)=\gcd(32,\ell)\in\{1,2,4\},
\]
and this gcd divides $5-1$.
Every prime factor of $N=\operatorname{lcm}(q,\lambda(q))$
divides $32m\ell$, so the residues in \eqref{bkapp:CRT} make $u$
a unit modulo $N$.
The second congruence fixes the values of $\beta$ and hence
$\beta^u=\beta$.  The first gives $u\equiv5\pmod8$, so
\[
 \alpha^u=\alpha^5=\alpha\eta_8.
\]
For the last identity it suffices to check the generators:
$\alpha^4(-1)=1$ and
$\alpha^4(5)=\zeta_8^{4j}=-1=\eta_8(5)$.
Thus $\chi^u=\chi\eta_8$ and
$\eta_8(u)=\eta_8(5)=-1$.
Lemma~\ref{bkapp:negative-swap} proves independence of the pair.

Now suppose that $j$ is even.  The kernel of
$U_{32}\longrightarrow U_{16}$ is $\{1,17\}$, and
$17=5^4\pmod{32}$.  Thus
\[
 \alpha(17)=\zeta_8^{4j}=1,
\]
so $\alpha$ factors modulo $16$.  The same is true of
$\alpha\eta_8$, whose exponent at $5$ is $j+4$.
The two distinct odd characters in the pair therefore descend to
distinct odd characters modulo $v=16m$.  Their arithmetic functions
on the integers are unchanged: the two moduli $32m$ and $16m$ have
the same prime divisors, so both the values on units and the zero
values on nonunits agree.  Their Dirichlet $L$-values are consequently
unchanged, without an additional Euler factor.
But
\[
 \lambda(v)=\operatorname{lcm}(4,\ell),
 \qquad \gcd(v,\lambda(v))=4,
\]
so the singleton-fibre case of Lemma~\ref{bkapp:fibres} proves
independence of all odd values modulo $v$, and in particular of this
pair.

Both cases are now settled.  Restriction-fibre splitting proves
independence of all odd values modulo $32m$.  The two established
parity facts stated above then imply independence of all
nonprincipal values.  Finally, for an odd prime $r$,
$\lambda(r)=r-1$ and $\gcd(r,r-1)=1$, so the displayed prime family
follows immediately.
\end{proof}

\begin{theorem}[Transcendence consequence]
\label{bkapp:transcendence}
For a modulus $q=32m$ satisfying the hypotheses of
Theorem~\ref{bkapp:main}, every rational linear combination
\[
 \sum_{\substack{\chi\bmod q\\\chi\ne\chi_0}}
   c_\chi L(1,\chi),
 \qquad c_\chi\in\mathbb Q,
\]
whose coefficients are not all zero is transcendental.
\end{theorem}

\begin{proof}
Theorem~\ref{bkapp:main} makes the displayed value nonzero.
For completeness, if $f$ is a periodic algebraic-valued function of
period $q$ with mean zero, Fourier inversion and Abel summation give
\[
 L(1,f)=-\sum_{a=1}^{q-1}\widehat f(a)\log(1-\zeta_q^a),
 \qquad
 \widehat f(a)=\frac1q\sum_{b=1}^q f(b)\zeta_q^{-ab},
\]
with the logarithms obtained by radial continuation from the unit
disc.  Apply this to $f=\sum_\chi c_\chi\chi$.
All coefficients and logarithm arguments are algebraic, and all
arguments are nonzero.  Baker's theorem on linear forms in
logarithms \cite{BK_Baker1975} says that such a linear form, if
nonzero, is transcendental.  This proves the assertion.
\end{proof}

The individual transcendence of the $L(1,\chi)$ is classical.
The consequence here is simultaneous nonvanishing, and hence
transcendence, for every nontrivial rational combination in the
stated family.  It does not assert independence of the odd values
over $\overline{\mathbb Q}$; those values are all algebraic
multiples of $\pi$.

\paragraph{Scope.}
The proof gives the full nonprincipal-character assertion for every
modulus in Theorem~\ref{bkapp:main}, rather than a rank estimate or a
finite verification.  It also proves that the ratios for pairs with primitive $32$-part are
algebraic and irrational.  It makes essential
use of two-element restriction fibres and the sign
$\eta_8(u)=-1$.  Neither the proof nor the theorem asserts the
corresponding result for arbitrary moduli or for every power of $2$.

\clearpage
\section{Reproducible finite illustrations}\label{app:checks}

The proofs above do not depend on a computation. The following standalone
Python~3 program illustrates Lemma~\ref{lem:values} by exhausting
$377\,304$ polynomials over the listed small fields, and checks the
finite-union warning of Remark~\ref{rem:union} at every odd prime below
$100$. It uses only the standard library. The expected output is
\texttt{PASS: 377304 polynomials; 24 odd primes}. These finite tests do
not prove the statements for arbitrary fields or arbitrary parameters,
and are not proof-assistant certificates.

\begingroup\small
\begin{verbatim}
from itertools import product
from math import isqrt


def is_prime(n):
    return n > 1 and all(n % d for d in range(2, isqrt(n) + 1))


count = 0
for p in (3, 5, 7, 11, 13):
    for d in range(2, 5):
        if (p - 1) % d:
            continue
        bound = p - (p - 1) // d
        powers = [[pow(x, j, p) for j in range(d + 1)]
                  for x in range(p)]
        for low in product(range(p), repeat=d):
            for lead in range(1, p):
                coefficients = low + (lead,)
                image = {
                    sum(c * v for c, v in zip(coefficients, row)) % p
                    for row in powers
                }
                assert len(image) <= bound, (p, d, coefficients)
                count += 1

prime_count = 0
for p in range(3, 100):
    if not is_prime(p):
        continue
    squares = {x * x % p for x in range(p)}
    for t in range(p):
        assert (t in squares or (t - 1) % p in squares
                or t * (t - 1) % p in squares), (p, t)
    prime_count += 1

assert count == 377304 and prime_count == 24
print(f"PASS: {count} polynomials; {prime_count} odd primes")
\end{verbatim}
\endgroup

\clearpage
\begin{thebibliography}{99}
\bibitem{Pollack}
P.~Pollack,
\emph{A simple proof of a theorem of Hajdu--Jarden--Narkiewicz},
Colloquium Mathematicum \textbf{147} (2017), 217--220.
\href{https://doi.org/10.4064/cm7054-9-2016}{doi:10.4064/cm7054-9-2016}.
Author manuscript: \url{https://www.pollack-math.net/sunit-impan.pdf}.

\bibitem{BL}
T.~Browning and D.~Loughran,
\emph{Sieving rational points on varieties},
Transactions of the American Mathematical Society \textbf{371} (2019),
no.~8, 5757--5785.
\href{https://doi.org/10.1090/tran/7514}{doi:10.1090/tran/7514}.
In particular, Lemma~3.8(2).

\bibitem{Serre}
J.-P.~Serre,
\emph{Topics in Galois Theory}, second edition,
Research Notes in Mathematics, vol.~1, A K Peters, 2008.
Chapter~3, on thin sets and Hilbert irreducibility.

\bibitem{CZ}
P.~Corvaja and U.~Zannier,
\emph{Rational and integral values of rational functions at rational points},
with an appendix by D.~Masser, preprint, 28 August 2026.
\href{https://arxiv.org/abs/2608.28255v1}{arXiv:2608.28255v1}.

\bibitem{Corpus60}
KEIO,
\emph{Exponent collapse and a question of Corvaja and Zannier on squares
in $S$-unit shifts}, research manuscript [060],
Verifiable Mathematical Reasoning Corpus,
repository snapshot \texttt{510a32c9a3fd2aa0b5c78d3f6812ded723b9ac8f},
consulted 28 October 2026.
Theorem~A treats the earlier even-degree polynomial setting over $\Q$.

\bibitem{Corpus53}
KEIO,
\emph{Restriction classes and Baker's question on the linear independence
of $L(1,\chi)$}, research manuscript [053],
Verifiable Mathematical Reasoning Corpus, dated 24 October 2026,
repository snapshot \texttt{510a32c9a3fd2aa0b5c78d3f6812ded723b9ac8f},
consulted 28 October 2026.

\bibitem{DZ}
P.~D\`ebes and U.~Zannier,
\emph{Universal Hilbert subsets},
Mathematical Proceedings of the Cambridge Philosophical Society
\textbf{124} (1998), 127--134.
\href{https://doi.org/10.1017/S0305004197001849}{doi:10.1017/S0305004197001849}.

\bibitem{Zannier}
U.~Zannier,
\emph{Diophantine equations with linear recurrences: An overview of some recent progress},
Journal de Th\'eorie des Nombres de Bordeaux \textbf{17} (2005), 423--435.
\url{https://www.numdam.org/article/JTNB_2005__17_1_423_0.pdf}.

\bibitem{BK_Baker1975}
A.~Baker, \emph{Transcendental Number Theory},
Cambridge University Press, Cambridge, 1975.

\bibitem{BK_BBW1973}
A.~Baker, B.~J.~Birch and E.~A.~Wirsing,
On a problem of Chowla,
\emph{J. Number Theory} \textbf{5} (1973), 224--236.

\bibitem{BK_MM2011}
M.~Ram Murty and V.~Kumar Murty,
A problem of Chowla revisited,
\emph{J. Number Theory} \textbf{131} (2011), no.~9, 1723--1733.
doi:10.1016/j.jnt.2011.04.003.

\bibitem{BK_KR2022}
N.~Kandhil and P.~Rath,
Around a question of Baker,
arXiv:2212.00358v1 (2022).
Published in \emph{Some Contributions to Number Theory and Beyond},
Fields Institute Communications \textbf{90}, Springer, 2026,
pp.~227--243, doi:10.1007/978-3-032-10837-1\_11.
The theorem numbering cited in the text is that of the arXiv version.

\bibitem{BK_MS2010}
M.~Ram Murty and N.~Saradha,
Euler--Lehmer constants and a conjecture of Erd\H{o}s,
\emph{J. Number Theory} \textbf{130} (2010), no.~12, 2671--2682.

\end{thebibliography}
\end{document}
